\documentclass[11pt]{article}

% Packages
\usepackage[margin=1in]{geometry}
\usepackage{amsmath,amssymb,amsfonts}
\usepackage{booktabs}
\usepackage{graphicx}
\usepackage{hyperref}
\usepackage{listings}
\usepackage{xcolor}
\usepackage{natbib}
\usepackage{enumitem}
% authblk removed, using standard \author
\usepackage[T1]{fontenc}
\usepackage[utf8]{inputenc}
\usepackage{CJKutf8}

% Listings style
\lstset{
  basicstyle=\ttfamily\small,
  breaklines=true,
  frame=single,
  language=Python,
  showstringspaces=false,
  columns=fullflexible,
  keepspaces=true,
  xleftmargin=1em,
  xrightmargin=1em,
}

% Hyperref setup
\hypersetup{
  colorlinks=true,
  linkcolor=blue,
  citecolor=blue,
  urlcolor=blue,
}

\title{vScore: Pre-Linguistic Visual Intelligence Through Domain-Agnostic Valence Scoring}

\author{Francois Reeves\thanks{Independent Researcher. B.A. Communications, McGill University. Data Science, Universit\'{e} de Paris 1 Sorbonne.} \and Research Assistant: Claude (Anthropic)}

\date{}

\begin{document}

\maketitle

\begin{abstract}
We propose vScore, a framework for visual intelligence that operates entirely below the level of language. Rather than mapping visual input to text labels or token sequences, vScore maps video directly to numerical valence vectors: continuous scores on domain-specific outcome axes where zero represents homeostasis and any deviation represents a cost demanding attention. The framework builds on the observation that biological intelligence, from insects to expert humans, evaluates visual stimuli through pre-linguistic pattern-to-outcome projection rather than word retrieval. We formalise this as a metaclass architecture where any scoreable domain (survival, fire, weather, sport, trading) is defined solely by its outcome axes, with a shared visual encoder providing domain-agnostic features. Using V-JEPA~2 as the frozen encoder and a set of lightweight valence heads, we demonstrate on 393 videos across 25 categories from three datasets (Kinetics-mini, HACS, THETIS) that (1)~a small regression head can learn to map self-supervised video features to valence scores with MAE of 0.79 within domain, (2)~three axes (coordination, impact, and speed) transfer across held-out domains with mean MAE of 1.45 to 1.74 (transferring in 62--85\% of holdout conditions), confirming the existence of universal visual primitives, (3)~axes with no visual analogue in the training set (e.g., synchronised coordination, collision impact) correctly fail to transfer, identifying genuinely novel visual dynamics, (4)~actions emerge from the \emph{geometry} of the full valence vector, specifically its direction rather than its magnitude, with trajectory projection enabling preemptive responses before threshold states are reached, and (5)~a Bayesian experience memory with posterior-aware eviction selectively retains informative experiences (high surprise, extreme outcomes, novel patterns) while discarding redundant ones, using normalised leave-one-out influence, coverage-driven eviction, and periodic replay to eliminate the sequential biases inherent in naive memory management. We argue that this valence-first, language-last hierarchy offers a more faithful model of biological visual intelligence than current vision-language paradigms. In this hierarchy, intelligence is geometric inference over outcome-projected vectors, learning is Bayesian posterior updating, memory is statistically rigorous retention, and words are an optional lookup table appended after scoring.
\end{abstract}

\section{Introduction}

The dominant paradigm in visual AI maps pixels to words. Image classifiers produce text labels. Vision-language models ground visual features in token sequences. Even self-supervised learners are typically evaluated by how well their representations support linguistic categorisation. The assumption is implicit: understanding an image means being able to describe it.

Biology disagrees. A gazelle seeing a predator does not produce the label ``lion'' before fleeing. A firefighter reading a blaze does not think ``the fire is spreading at rate X toward exit Y'' before acting. A hockey coach does not narrate the play before calling a timeout. In each case, the visual stimulus is evaluated on outcome-relevant axes (threat, speed, containment, momentum) and the response is triggered by threshold crossing on those axes. Language, when it arrives, is a post-hoc serialisation of a decision already made.

This intuition, that intelligence requires world models operating in representation space rather than token space, was articulated by \citet{lecun2022} in his vision for autonomous machine intelligence. LeCun argued that autoregressive language models are fundamentally limited because they operate on discrete tokens rather than continuous representations of the world, and proposed the Joint Embedding Predictive Architecture (JEPA) as the foundation for a system that learns by predicting in latent space rather than pixel or token space. V-JEPA and V-JEPA~2, the self-supervised video encoders we build upon, are direct implementations of this vision. vScore extends LeCun's framework in a specific direction: if the encoder learns to predict the world in latent space without language, what should the downstream evaluation look like? Our answer is that it should look like biological valence scoring: not classification, not captioning, but continuous outcome projection on domain-relevant axes.

We chose the term \textit{valence} deliberately. In chemistry, valence describes an atom's capacity to bond along specific axes, with electrons orbiting a stable nucleus. In our framework, the valence vector describes an organism's capacity to \textit{react} along domain-specific axes, orbiting a stable zero point (homeostasis). The analogy is structural, not metaphorical: just as an electron's orbital shape is defined by its quantum numbers, a valence vector's meaning is defined by the axes it deviates along. The name vScore reflects this: vectors of scored reactions, orbiting expected equilibrium in domain-specific experience space.

This paper introduces vScore, a framework that inverts the standard hierarchy. Instead of:

\begin{lstlisting}[language={}]
pixels -> encoder -> language -> understanding
\end{lstlisting}

\noindent we propose:

\begin{lstlisting}[language={}]
pixels -> encoder -> valence scores -> trajectory projection -> threshold trigger -> [language, optionally]
\end{lstlisting}

The key contributions are:

\begin{enumerate}
    \item \textbf{A metaclass formalisation} where any domain with scoreable outcomes (survival, fire, sport, trading) is defined by its numerical axes alone, enabling domain-agnostic architecture with domain-specific scoring.

    \item \textbf{Empirical evidence} that self-supervised video features (V-JEPA~2) contain universal visual primitives, notably motion dynamics, that transfer across visually distinct domains without any language supervision.

    \item \textbf{A threshold-triggering mechanism} with dynamic thresholds that lower under acceleration, modelling the biological principle that organisms act on projected trajectories, not current states.

    \item \textbf{A hierarchy} in which language is the final, optional layer: a multilingual lookup table indexed by valence state, not a prerequisite for understanding.
\end{enumerate}

\section{Motivation: The Pre-Linguistic Layer}

\subsection{The Translation Problem}

Every word is already a translation. When a human says ``fire,'' they have already performed a complex chain of operations: detect high-contrast flickering pattern, evaluate colour (orange-red), estimate spatial extent, assess motion (spreading vs.\ contained), project trajectory (toward me or away), integrate with context (indoor vs.\ outdoor, near exits or not), cross a relevance threshold, retrieve a lexical item from memory, and articulate. The word ``fire'' compresses all of this into four letters. A vision-language model trained on the word ``fire'' learns the compression, not the evaluation chain that produced it.

This matters because the evaluation chain is where intelligence lives. Two fires can both be labelled ``fire'' while demanding opposite responses. One is a campfire (approach, warmth, safety) and the other is a structure fire (flee, danger, urgency). The word is the same. The valence is opposite.

\subsection{Biological Precedent}

\citet{panksepp1998} identified seven primal affective circuits operating at the subcortical level in all mammals: SEEKING, RAGE, FEAR, LUST, CARE, PANIC/GRIEF, and PLAY. These circuits:

\begin{itemize}
    \item Operate before and independently of cortical (linguistic) processing
    \item Are activated by specific sensory patterns, not labels
    \item Produce graded responses (not binary classifications)
    \item Interact with each other (FEAR suppresses PLAY; SEEKING modulates all others)
    \item Drive behaviour through threshold dynamics, not categorisation
\end{itemize}

\citeauthor{gibson1979}'s \citeyearpar{gibson1979} theory of affordances makes a complementary argument: perception is not about identifying objects but about detecting action possibilities. A surface is not perceived as ``a chair'' but as ``sittable.'' The affordance is pre-linguistic and directly coupled to the organism's body and goals.

\citeauthor{damasio1994}'s \citeyearpar{damasio1994} somatic marker hypothesis extends this further: decisions are driven by body-state projections, essentially emotional previews of future outcomes, not by rational analysis. The ``gut feeling'' is not noise; it is the organism's fastest evaluation system.

vScore formalises these biological principles computationally.

\subsection{The Expert's Eye}

The gap between novice and expert in any visual domain is not vocabulary but scoring speed and accuracy. A rookie firefighter and a 20-year veteran both know the word ``flashover.'' The difference is that the veteran sees the visual precursors (darkening smoke, thermal layering, rollover at the ceiling) and scores the trajectory toward flashover before it happens. The expert's advantage is entirely pre-linguistic: faster feature extraction, more accurate valence scoring, better trajectory projection.

This suggests that expert training is fundamentally about calibrating valence heads, not expanding vocabulary. The vScore framework makes this explicit.

\section{Framework}

\subsection{Architecture Overview}

vScore consists of four components arranged in a strict hierarchy:

\begin{lstlisting}[language={}]
Level 2: Visual Encoder (frozen, self-supervised)
    Input: Video frames (T, C, H, W)
    Output: Dense feature tensor (N_tokens, D)

Level 1: Valence Head (trainable, per-domain)
    Input: Pooled features (D,)
    Output: Valence vector (N_axes,)

Level 0: Trajectory Projector + Threshold Trigger
    Input: Sequence of valence vectors over time
    Output: Trigger events (IGNORE / ATTEND / ACT)

Level 3 (optional): Language Lookup
    Input: Valence state
    Output: Multilingual text description
\end{lstlisting}

Levels 0--2 are the intelligence. Level~3 is serialisation.

\subsection{The ScoredDomain Metaclass}

We formalise domain definition using a Python metaclass:

\begin{lstlisting}[language=Python]
class ScoredDomain(type):
    registry = {}
    def __new__(mcs, name, bases, namespace):
        cls = super().__new__(mcs, name, bases, namespace)
        if 'axes' in namespace and namespace['axes'] is not None:
            cls._axis_index = {ax: i for i, ax in enumerate(axes)}
            cls._n_axes = len(axes)
            mcs.registry[name] = cls
        return cls
\end{lstlisting}

A domain is defined entirely by its axes:

\begin{lstlisting}[language=Python]
class Survival(metaclass=ScoredDomain):
    axes = ['seeking', 'rage', 'fear', 'lust', 'care', 'panic', 'play']

class Fire(metaclass=ScoredDomain):
    axes = ['spread_rate', 'proximity', 'intensity', 'containment',
            'escape_route', 'structural_risk', 'smoke_toxicity']

class Hockey(metaclass=ScoredDomain):
    axes = ['scoring_threat', 'possession_pressure', 'breakaway',
            'penalty_risk', 'momentum', 'fatigue']
\end{lstlisting}

This design enforces a critical constraint: \textbf{the model architecture is identical across domains}. Only the number and meaning of axes change. A fire valence head and a hockey valence head are the same neural network with different output dimensions. The visual encoder is shared.

\subsection{Valence Vectors and Homeostasis}

A valence vector $\mathbf{v} \in \mathbb{R}_{\geq 0}^{n}$ scores the current state on $n$ domain-specific axes. The constraint $v_i \geq 0$ is enforced architecturally (ReLU activation on the output layer).

\textbf{Zero is homeostasis.} The organism at rest, no activation needed. Every non-zero value represents a deviation that demands energy, whether attention, computation, or physical action. The system's purpose is to return to zero.

This is a fundamentally different framing from classification. A classifier asks ``what is this?'' A valence scorer asks ``how much does this deviate from nothing happening, and in which directions?''

The magnitude $\|\mathbf{v}\|_2$ gives overall activation level. A video of an empty field scores near zero. A video of a building collapse scores high on multiple axes simultaneously.

\subsection{Trajectory Projection}

A single valence vector is a snapshot. Intelligence requires trajectory: where is each axis heading?

Given a sequence of valence vectors $\mathbf{v}_{t-k}, \ldots, \mathbf{v}_t$, we compute per-axis:

\begin{itemize}
    \item \textbf{Velocity}: $\dot{v}_i = \frac{v_i^{(t)} - v_i^{(t-k)}}{k \cdot \Delta t}$, the rate of change
    \item \textbf{Acceleration}: $\ddot{v}_i$, the rate of change of rate of change
    \item \textbf{Projection}: $\hat{v}_i^{(t+h)} = v_i^{(t)} + \dot{v}_i \cdot h$, the expected axis value at horizon $h$
\end{itemize}

The linear projection serves as baseline. The learned TemporalProjector, a small transformer over valence sequences, captures nonlinear dynamics such as the super-linear acceleration of fire spread following containment loss.

\subsection{Dynamic Threshold Triggering}

Biological systems do not classify; they trigger. The response function is:

\begin{equation}
\text{response}(v_i, \hat{v}_i) = \begin{cases} \text{ACT} & \text{if } \max(v_i, \hat{v}_i) \geq \tau_i - \alpha \cdot \max(0, \ddot{v}_i) \\ \text{ATTEND} & \text{if } \max(v_i, \hat{v}_i) \geq \tau_i^{\text{att}} - \frac{\alpha}{2} \cdot \max(0, \ddot{v}_i) \\ \text{IGNORE} & \text{otherwise} \end{cases}
\end{equation}

\noindent where $\tau_i$ is the base trigger threshold, $\tau_i^{\text{att}}$ is the attention threshold, and $\alpha$ is the acceleration weight.

The key mechanism: \textbf{positive acceleration lowers the threshold.} A fire spreading at constant speed triggers at FEAR=7.0. A fire that is accelerating triggers at FEAR=5.5. The organism does not wait for peak danger; it acts on the derivative.

This models a well-documented biological phenomenon: startle responses and preemptive action are driven by rate of change, not absolute magnitude \citep{schiff1965,regan1995}.

\subsection{Action Inference from Valence Geometry}

Per-axis thresholds (Section~3.5) answer ``is this axis critical?'' But real decisions emerge from the \textbf{interaction} between axes, from the geometry of the full valence vector. A fire with \texttt{[spread=7, containment=8]} is a controlled burn. The same spread with \texttt{[spread=7, containment=2]} is a crisis. Same axis, same score, opposite action. The difference is the vector's direction in valence space.

We formalise this through an ActionField: a mapping from regions of valence space to response tendencies. Each ActionRegion is defined by a condition function over the full vector, not individual axes:

\begin{equation}
a_k(\mathbf{v}) = f_k(v_1, v_2, \ldots, v_n) \in [0, 1]
\end{equation}

\noindent where $a_k$ is the activation of action $k$ and $f_k$ captures the interaction logic. For example, in the survival domain:

\begin{table}[h]
\centering
\begin{tabular}{lll}
\toprule
\textbf{Action} & \textbf{Condition} & \textbf{Biological basis} \\
\midrule
FLEE & fear dominant, rage low & Predator escape \\
FIGHT & rage dominant, fear low & Territory defence \\
FREEZE & fear AND rage both high & Competing drives, tonic immobility \\
PROTECT & care AND fear both high & Parental defence despite danger \\
EXPLORE & seeking high, fear low & Foraging, curiosity \\
\bottomrule
\end{tabular}
\end{table}

Three properties make this system qualitatively different from classification:

\textbf{1. Direction, not magnitude.} Four vectors with identical L2 norm ($\sim$10.0) produce four completely different actions:

\begin{table}[h]
\centering
\small
\begin{tabular}{llll}
\toprule
\textbf{Vector} & \textbf{Action} & \textbf{Activation} & \textbf{Conflict} \\
\midrule
\texttt{[0,0,10,0,0,0,0]} (pure fear) & FLEE (1.00) & 1.00 & 0.33 \\
\texttt{[0,10,0,0,0,0,0]} (pure rage) & FIGHT (1.00) & 1.00 & 0.33 \\
\texttt{[0,0,0,0,10,0,0]} (pure care) & BOND (1.00) & 1.00 & 0.33 \\
\texttt{[3.8,3.8,3.8,3.8,3.8,3.8,3.8]} (distributed) & FREEZE (0.38) & 0.38 & 0.53 \\
\bottomrule
\end{tabular}
\end{table}

A classifier would label all four as ``high activation.'' The valence scorer identifies four qualitatively different states. Magnitude is arousal. Direction is meaning.

\textbf{2. Conflict detection.} When multiple action regions activate at similar strength, the system quantifies conflict as the ratio of the second-strongest to the strongest activation:

\begin{equation}
\text{conflict}(\mathbf{v}) = \frac{a_{\text{second}}(\mathbf{v})}{a_{\text{top}}(\mathbf{v})}
\end{equation}

Conflict near 1.0 means the organism faces equally compelling but incompatible actions, the classical Buridan's donkey problem. This is biologically real: tonic immobility in prey animals \citep{gallup1977}, decision paralysis in humans under competing threat and reward signals \citep{corr2013}. The conflict score predicts hesitation latency, not just action identity.

In the fire domain, we observe high conflict ($>$0.7) at the transition point where ADVANCE and HOLD\_POSITION compete, exactly when a firefighter's decision is hardest. The model quantifies what experienced commanders simply call ``the judgement call.''

\textbf{3. Preemptive action from trajectory.} The action field evaluates not only the current valence vector but the projected vector at horizon $h$. When the projected state activates a higher-priority action than the current state, the system triggers a preemption signal:

\begin{equation}
\text{preempt} = \begin{cases} \text{true} & \text{if } \text{top}_k(a(\hat{\mathbf{v}}_{t+h})) \neq \text{top}_k(a(\mathbf{v}_t)) \text{ and } a_{\text{top}}(\hat{\mathbf{v}}_{t+h}) > a_{\text{top}}(\mathbf{v}_t) \\ \text{false} & \text{otherwise} \end{cases}
\end{equation}

In our fire simulation, preemption fires at $t=2$ (current action: HOLD\_POSITION; projected action: DEFENSIVE) and at $t=4$ (current: EVACUATE; projected: MAYDAY). The message is clear: \emph{the present is manageable but the future is not. Act on the projection.} This is how expert firefighters survive: they read the trajectory and leave before the current state demands it.

The combination of direction-dependent action, conflict quantification, and trajectory-based preemption produces behaviour that is qualitatively richer than threshold-crossing on individual axes. A 7-dimensional valence space with trajectory generates a continuous manifold of behavioural states, each point defined not by which axis is highest, but by the angle of the vector, the curvature of its trajectory, and the proximity of competing action regions.

\subsection{The Language Layer (Optional)}

Language enters only at Level~3, as a lookup from valence state to multilingual text:

\begin{lstlisting}[language={}]
[0, 0, 8, 0, 0, 3, 0] -> FEAR-dominant, PANIC-secondary
    EN: "charging bear"
    FR: "ours qui charge"
    JP: "tosshinsuru kuma" (charging bear)
    Physiological: HR up, pupil dilation, cortisol spike
\end{lstlisting}

The text is not the understanding. The understanding is the valence vector. The text is one possible serialisation, chosen for a specific audience in a specific language. The same understanding could be serialised as a motor command (flee), a physiological response (adrenaline), or a social signal (warning call). All are downstream of the same valence evaluation.

\subsection{Dual-System Integration: vScore + LLM}

vScore does not replace language models. It sits before them. The two systems are complementary, operating at different speeds on different representations:

\begin{table}[h]
\centering
\small
\begin{tabular}{lll}
\toprule
\textbf{Property} & \textbf{System 1 (vScore)} & \textbf{System 2 (LLM)} \\
\midrule
Speed & $\sim$50ms & $\sim$500ms+ \\
Input & Patterns (pixels, waveforms, prosody) & Tokens (words, syntax) \\
Output & Valence vectors, action triggers & Semantic interpretation, reasoning \\
Strength & Speed, cross-modal fusion, threat detection & Context, planning, social reasoning \\
Weakness & No reasoning, no planning & Slow, requires language, misses prosody \\
\bottomrule
\end{tabular}
\end{table}

The critical architectural principle: \textbf{vScore gates the LLM, not the reverse.} The fast system decides whether the slow system is needed:

\begin{itemize}
    \item \textbf{Low valence magnitude} ($\|\mathbf{v}\| < \tau_{\text{floor}}$): nothing happening, neither system activates
    \item \textbf{High valence, low conflict}: clear situation, vScore acts alone (fleeing from a charging bear leaves no time for words)
    \item \textbf{High valence, high conflict}: ambiguous situation, engage the LLM for reasoning (a child laughing then crying activates both CARE and PANIC, and context is needed)
    \item \textbf{Prosody/semantics mismatch}: voice intonation (scored by vScore) contradicts word content (parsed by LLM), suggesting possible deception, sarcasm, or masked distress
\end{itemize}

The prosodic bridge is particularly significant. Voice intonation is an acoustic pattern with pitch contour, tempo, roughness, and rhythm, all scoreable on valence axes without understanding a single word. vScore processes prosody as sound: a raised voice activates threat regardless of the language spoken. The LLM processes the words. When prosodic valence and semantic content diverge (for example, ``I'm fine'' spoken with a trembling voice) the system detects the mismatch and flags it for deliberative processing.

This mirrors the known architecture of biological threat processing: the amygdala (fast, subcortical) evaluates sensory input and gates prefrontal cortex (slow, deliberative) engagement. The startle response precedes conscious evaluation. ``I flinched before I knew why'' is vScore firing before the LLM started.

\section{Experiments}

\subsection{Setup}

\textbf{Encoder.} We use V-JEPA~2 (ViT-L, 0.3B parameters) pretrained via self-supervised video prediction \citep{bardes2025}. The model is frozen throughout all experiments. It produces 8{,}192 spatiotemporal tokens of dimension 1{,}024 per 64-frame video clip. We global-average-pool to a single 1{,}024-dimensional vector per video.

\textbf{Data.} 393 videos from three sources: Kinetics-mini (150 videos, train+val splits, 5 categories: archery, bowling, flying kite, high jump, marching), HACS \citep{zhao2019} (63 videos, validation segments, 8 categories: arm wrestling, fencing, fixing roof, javelin throw, pole vault, slacklining, snowboarding, springboard diving), and THETIS \citep{gourgari2013} (180 videos, 12 tennis shot types from 55 players: backhand, backhand two-handed, backhand slice, backhand volley, forehand flat, forehand open stance, forehand slice, forehand volley, flat serve, kick serve, slice serve, smash). Each clip is sampled at 64 frames uniformly. Features are extracted once and cached as 1{,}024-dimensional vectors; all training experiments operate on cached vectors.

\textbf{Domain.} We define a ``dynamics'' domain with 6 axes capturing universal motion properties:

\begin{table}[h]
\centering
\begin{tabular}{ll}
\toprule
\textbf{Axis} & \textbf{Description} \\
\midrule
speed & Overall motion magnitude \\
impact & Collision / force transfer \\
precision & Controlled vs.\ chaotic motion \\
verticality & Up/down movement dominance \\
coordination & Synchronised multi-agent motion \\
tension & Buildup before release/outcome \\
\bottomrule
\end{tabular}
\end{table}

\textbf{Proxy annotations.} Each category receives a fixed valence vector based on its dominant visual dynamics (Table~\ref{tab:proxy}). This is a deliberate simplification; real deployment would use per-video expert annotations or physiological measurements.

\begin{table}[t]
\centering
\caption{Proxy valence scores (13 categories, 2 datasets).}
\label{tab:proxy}
\small
\begin{tabular}{llccccccc}
\toprule
\textbf{Category} & \textbf{Source} & \textbf{N} & \textbf{speed} & \textbf{impact} & \textbf{precision} & \textbf{verticality} & \textbf{coordination} & \textbf{tension} \\
\midrule
archery & K-mini & 30 & 3.0 & 2.0 & 9.0 & 1.0 & 1.0 & 8.0 \\
bowling & K-mini & 30 & 5.0 & 8.0 & 7.0 & 0.5 & 1.0 & 6.0 \\
flying kite & K-mini & 30 & 4.0 & 0.5 & 3.0 & 7.0 & 1.0 & 1.0 \\
high jump & K-mini & 30 & 7.0 & 4.0 & 6.0 & 9.0 & 1.0 & 7.0 \\
marching & K-mini & 30 & 3.0 & 1.0 & 5.0 & 0.5 & 9.0 & 1.0 \\
arm wrestling & HACS & 3 & 2.0 & 6.0 & 5.0 & 0.5 & 2.0 & 9.0 \\
fencing & HACS & 7 & 7.0 & 4.0 & 9.0 & 1.0 & 2.0 & 8.0 \\
fixing roof & HACS & 9 & 2.0 & 3.0 & 5.0 & 5.0 & 1.0 & 4.0 \\
javelin throw & HACS & 10 & 8.0 & 3.0 & 8.0 & 6.0 & 1.0 & 8.0 \\
pole vault & HACS & 10 & 7.0 & 3.0 & 8.0 & 9.0 & 1.0 & 9.0 \\
slacklining & HACS & 4 & 2.0 & 1.0 & 8.0 & 3.0 & 1.0 & 7.0 \\
snowboarding & HACS & 10 & 8.0 & 4.0 & 6.0 & 4.0 & 1.0 & 4.0 \\
spr.\ diving & HACS & 10 & 5.0 & 5.0 & 8.0 & 9.0 & 1.0 & 7.0 \\
\bottomrule
\end{tabular}
\end{table}

\textbf{Valence head.} A 3-layer MLP ($1024 \to 256 \to 128 \to 6$) with GELU activations and ReLU output. Trained with MSE loss, Adam optimiser (lr=$1\mathrm{e}{-3}$), 500 epochs.

\subsection{Cross-Domain Feature Analysis}

Before training valence heads, we analyse the raw V-JEPA~2 feature space across three representative categories (archery, bowling, flying kite).

\begin{table}[h]
\centering
\caption{Global cosine similarity between domain-pooled features.}
\label{tab:cosine}
\begin{tabular}{lccc}
\toprule
 & \textbf{archery} & \textbf{bowling} & \textbf{flying kite} \\
\midrule
archery & 1.00 & 0.45 & 0.69 \\
bowling & 0.45 & 1.00 & 0.56 \\
flying kite & 0.69 & 0.56 & 1.00 \\
\bottomrule
\end{tabular}
\end{table}

The encoder produces features that are neither identical across domains (which would prevent discrimination) nor orthogonal (which would prevent transfer). The intermediate similarity (0.45--0.69) suggests a shared visual vocabulary with domain-specific variation, precisely the structure needed for domain-agnostic scoring with domain-specific heads.

Token-level analysis reveals 8{,}192 spatiotemporal tokens per video with mean self-similarity of 0.24--0.28 within domains and 0.12--0.18 across domains. Critically, maximum cross-domain token similarity reaches 0.54, meaning specific visual tokens (likely encoding motion dynamics) activate nearly identically across different domains.

We identify 20 feature dimensions with high minimum activation across all domains (universal primitives, likely encoding edges, motion magnitude, and spatial structure) and 20 dimensions with high variance across domains (domain-discriminating features encoding texture and scene-specific patterns).

\subsection{Within-Domain Training}

With random 80/20 train/test split across all 213 videos (13 categories):

\textbf{Overall test MAE: 0.79} (on a 0--10 scale)

\begin{table}[h]
\centering
\caption{Per-axis MAE (random split).}
\label{tab:peraxis}
\begin{tabular}{cccccc}
\toprule
\textbf{speed} & \textbf{impact} & \textbf{precision} & \textbf{verticality} & \textbf{coordination} & \textbf{tension} \\
\midrule
0.34 & 0.43 & 0.43 & 0.66 & 2.33 & 0.55 \\
\bottomrule
\end{tabular}
\end{table}

Five of six axes achieve MAE below 0.7, confirming that the valence head reliably maps V-JEPA~2 features to outcome-relevant scores. Coordination remains the hardest axis (2.33), consistent with the observation that synchronised multi-body motion is a higher-order visual pattern requiring more diverse training examples. The overall MAE of 0.79 represents a 36\% improvement over the initial 50-video experiment (1.24), demonstrating that the framework scales efficiently with data.

\subsection{Cross-Domain Transfer}

The critical experiment: train on 12 categories, test on the held-out 13th. If the model learns visual dynamics rather than category identity, axes should transfer to categories never seen during training.

\begin{table}[t]
\centering
\caption{Per-axis MAE when tested on held-out category (13 categories, 213 videos). Bold values indicate axes where transfer succeeds (MAE $< 2.0$).}
\label{tab:transfer}
\small
\begin{tabular}{lcccccccr}
\toprule
\textbf{Held-out} & \textbf{N} & \textbf{speed} & \textbf{impact} & \textbf{precision} & \textbf{verticality} & \textbf{coordination} & \textbf{tension} & \textbf{Overall} \\
\midrule
arm wrestling & 3 & 2.03 & 2.29 & \textbf{1.54} & \textbf{0.78} & \textbf{0.87} & 3.80 & 1.89 \\
fencing & 7 & \textbf{0.93} & \textbf{0.13} & 2.33 & 3.89 & \textbf{0.33} & 2.09 & 1.62 \\
fixing roof & 9 & 2.37 & \textbf{0.56} & 3.31 & 3.18 & \textbf{0.32} & 3.44 & 2.20 \\
javelin throw & 10 & \textbf{1.30} & \textbf{0.19} & \textbf{0.94} & \textbf{1.31} & \textbf{0.66} & \textbf{0.67} & 0.85 \\
pole vault & 10 & \textbf{0.69} & \textbf{0.31} & 2.52 & \textbf{0.43} & \textbf{0.22} & 9.00 & 2.20 \\
slacklining & 4 & 2.25 & \textbf{1.66} & \textbf{1.10} & \textbf{1.09} & \textbf{0.22} & \textbf{1.12} & 1.24 \\
snowboarding & 10 & 4.16 & 2.26 & 6.00 & \textbf{1.13} & \textbf{0.37} & \textbf{1.06} & 2.49 \\
spr.\ diving & 10 & 2.61 & \textbf{0.32} & \textbf{1.43} & \textbf{0.90} & \textbf{0.51} & \textbf{0.67} & 1.07 \\
archery & 30 & \textbf{0.82} & \textbf{0.73} & 4.31 & 3.20 & \textbf{1.70} & 5.10 & 2.64 \\
bowling & 30 & \textbf{1.99} & 5.35 & 2.51 & 2.18 & 3.70 & 3.14 & 3.15 \\
flying kite & 30 & \textbf{0.74} & \textbf{1.89} & 4.37 & 5.18 & \textbf{1.62} & 4.81 & 3.10 \\
high jump & 30 & \textbf{0.88} & \textbf{0.92} & \textbf{1.56} & \textbf{1.81} & \textbf{0.50} & \textbf{1.05} & 1.12 \\
marching & 30 & \textbf{1.90} & 2.49 & \textbf{1.33} & 2.99 & 7.86 & 4.70 & 3.54 \\
\bottomrule
\end{tabular}
\end{table}

\begin{table}[h]
\centering
\caption{Universality ranking (mean transfer MAE across all 13 holdouts).}
\label{tab:universality}
\begin{tabular}{lccc}
\toprule
\textbf{Axis} & \textbf{Mean MAE} & \textbf{Transfers in} & \textbf{Classification} \\
\midrule
coordination & 1.45 & 11/13 (85\%) & \textbf{Universal} \\
impact & 1.47 & 9/13 (69\%) & \textbf{Universal} \\
speed & 1.74 & 8/13 (62\%) & \textbf{Universal} \\
verticality & 2.16 & 7/13 (54\%) & Semi-universal \\
precision & 2.56 & 6/13 (46\%) & Semi-universal \\
tension & 3.13 & 5/13 (38\%) & Domain-specific \\
\bottomrule
\end{tabular}
\end{table}

\textbf{Key findings:}

\begin{enumerate}
    \item \textbf{Three axes are universal visual primitives.} Coordination (MAE 1.45, transfers in 85\% of holdouts), impact (1.47, 69\%), and speed (1.74, 62\%) generalise across visually distinct domains. The V-JEPA~2 encoder captures these motion dynamics in a form that is separable from scene identity and transferable to unseen categories.

    \item \textbf{Category diversity reveals hidden universality.} Coordination was classified as domain-specific in the 5-category experiment (MAE 8.09 when marching was the only source). With 13 categories including fencing, diving, slacklining, and pole vault, all involving body coordination, the axis transfers to 11/13 holdouts. The visual primitive was always in the encoder; it needed diverse training examples to be extracted.

    \item \textbf{Some categories transfer on all axes.} Javelin throw (MAE 0.85) and high jump (MAE 1.12) achieve successful transfer on all 6 axes. The remaining 12 categories collectively teach everything about their visual dynamics. This is the strongest evidence that universal visual primitives exist and are sufficient for scoring unseen domains.

    \item \textbf{Tension remains domain-specific.} Even with 13 categories, tension transfers in only 5/13 holdouts (MAE 3.13). The visual signature of ``buildup before release'' differs too much across archery, pole vault, and arm wrestling to generalise without domain-specific training. Notably, pole vault tension transfer fails catastrophically (MAE 9.00) because the stillness-before-the-run is visually unique.

    \item \textbf{Failure to transfer identifies novel dynamics.} Bowling's impact (5.35), snowboarding's speed (4.16), and marching's coordination (7.86) remain high-MAE holdouts, confirming that these categories contain visual patterns not represented elsewhere in the dataset. The system correctly identifies which domains need dedicated training data.
\end{enumerate}

\subsubsection{Random Baseline Control}

A natural objection is that 300M encoder parameters and a 400K-parameter MLP have enough capacity to produce low transfer MAE by chance, without learning anything about visual dynamics. To rule this out, we scramble the proxy scores: each category receives a random valence vector (uniform on $[0, 10]^6$), and the full cross-domain transfer experiment is repeated. We average over 5 random assignments.

\begin{table}[ht]
\centering
\caption{Real vs.\ random baseline transfer MAE. Ratio $< 0.7$ indicates strong signal.}
\begin{tabular}{lrrrl}
\toprule
Axis & Real MAE & Random MAE & Ratio & Verdict \\
\midrule
coordination & 1.68 & 3.18 & 0.53 & \textbf{SIGNAL} \\
speed & 1.73 & 2.81 & 0.62 & \textbf{SIGNAL} \\
impact & 1.86 & 3.01 & 0.62 & \textbf{SIGNAL} \\
precision & 2.14 & 3.33 & 0.64 & \textbf{SIGNAL} \\
verticality & 2.06 & 3.04 & 0.68 & \textbf{SIGNAL} \\
tension & 2.89 & 3.06 & 0.95 & noise \\
\midrule
Overall & 2.06 & 3.07 & 0.67 & \\
\bottomrule
\end{tabular}
\label{tab:random_baseline}
\end{table}

Five of six axes transfer 30--47\% better than random chance. Coordination shows the strongest signal (ratio 0.53): the model transfers nearly twice as well as a random assignment. Tension, at ratio 0.95, is indistinguishable from noise, confirming that it is genuinely domain-specific rather than merely undertrained.

The overall ratio of 0.67 rules out the capacity objection. If transfer were an artefact of parameter count or encoder clustering, random scores would transfer equally well. They do not. The valence head learns a structured mapping from visual features to outcome-relevant scores that random assignments cannot replicate.

\subsubsection{Single-Domain Validation: Tennis}

To test whether vScore operates within a single sport domain, we apply the framework to the THETIS tennis dataset (1,980 RGB videos, 55 players, 12 shot types). We extract V-JEPA~2 features from 15 videos per shot type (180 total) and define six tennis-specific axes: speed, impact, precision, verticality, aggression, and tension. Each shot type receives a proxy valence vector reflecting its dominant visual dynamics (e.g., flat serve: speed=9, tension=8; backhand slice: precision=8, speed=4).

\begin{table}[ht]
\centering
\caption{Cross-shot transfer on THETIS (train on 11 shot types, test on held-out 12th)}
\begin{tabular}{lrrrrrrrrr}
\toprule
Held-out & N & MAE & speed & impact & prec. & vert. & aggr. & tension \\
\midrule
backhand & 15 & \textbf{1.04} & 0.88* & 1.22* & 0.95* & 0.89* & 1.13* & 1.15* \\
backhand 2H & 15 & 1.33 & 1.45* & 2.31 & 1.10* & 0.82* & 1.74* & 0.54* \\
backhand slice & 15 & 2.13 & 1.34* & 2.11 & 0.63* & 2.00 & 2.68 & 4.00 \\
backhand volley & 15 & 2.02 & 0.69* & 0.88* & 8.00 & 0.53* & 1.15* & 0.90* \\
flat serve & 15 & 3.40 & 1.98* & 1.34* & 7.00 & 0.83* & 1.24* & 8.00 \\
forehand flat & 15 & \textbf{1.29} & 1.90* & 1.97* & 0.86* & 0.44* & 1.87* & 0.69* \\
forehand open & 15 & \textbf{1.30} & 1.57* & 1.81* & 0.69* & 1.31* & 1.46* & 0.94* \\
forehand slice & 15 & 1.42 & 1.16* & 2.22 & 0.80* & 0.23* & 2.63 & 1.45* \\
forehand volley & 15 & 2.15 & 0.86* & 1.36* & 0.85* & 0.91* & 7.00 & 1.93* \\
kick serve & 15 & 2.09 & 7.00 & 0.96* & 0.57* & 1.54* & 0.87* & 1.58* \\
slice serve & 15 & 2.04 & 2.14 & 2.29 & 1.69* & 2.45 & 1.93* & 1.72* \\
smash & 15 & 3.28 & 2.82 & 3.86 & 2.70 & 3.77 & 2.64 & 3.92 \\
\bottomrule
\end{tabular}
\label{tab:tennis_transfer}
\end{table}

\noindent * indicates MAE $< 2.0$ (successful transfer).

Three findings emerge from this single-domain experiment:

\begin{enumerate}
    \item \textbf{Visually similar shots are interchangeable.} Backhand (MAE 1.04), forehand flat (1.29), and forehand open stance (1.30) transfer on all six axes. The 11 remaining shot types collectively teach everything about these groundstrokes. A coach watching the same videos would agree: a forehand flat and a forehand open stance share body mechanics, racquet path, and contact dynamics. The encoder captures this without being told.

    \item \textbf{Visually unique shots correctly fail.} Smash (MAE 3.28) fails on all axes. Its overhead motion, extreme verticality, and downward strike have no analogue in groundstrokes or serves. Flat serve tension (MAE 8.00) fails because the service toss and wind-up are visually unlike any other shot's preparation phase. These failures are informative: they identify which shot dynamics require dedicated training.

    \item \textbf{Global pooling is the bottleneck.} The raw feature clustering analysis shows cosine similarity of 0.975--0.997 between all shot pairs, indicating that the global average over 8,192 tokens is dominated by the shared indoor background rather than the discriminative motion. The valence head partially overcomes this by learning which feature dimensions carry shot-specific information, but per-token or attention-weighted pooling would likely improve both clustering and transfer.
\end{enumerate}

The tennis experiment confirms that vScore generalises to a single sport domain with fine-grained action types. The framework requires no modification: the same metaclass, the same encoder, the same valence head architecture. Only the axes and their scores change.

\subsection{Action Inference from Valence Geometry}

To test whether multi-axis valence vectors produce qualitatively richer behaviour than per-axis thresholds, we construct action fields for the survival and fire domains and evaluate them on simulated valence trajectories.

\subsubsection{Direction vs.\ Magnitude (Survival)}

We construct four valence vectors with identical L2 magnitude ($\sim$10.0) but different directions in the 7-dimensional Panksepp space:

\begin{table}[h]
\centering
\caption{Same magnitude, different actions.}
\label{tab:direction}
\begin{tabular}{lcccc}
\toprule
\textbf{Vector profile} & \textbf{Magnitude} & \textbf{Top action} & \textbf{Activation} & \textbf{Conflict} \\
\midrule
Pure fear \texttt{[0,0,10,0,0,0,0]} & 10.0 & FLEE & 1.00 & 0.33 \\
Pure rage \texttt{[0,10,0,0,0,0,0]} & 10.0 & FIGHT & 1.00 & 0.33 \\
Pure care \texttt{[0,0,0,0,10,0,0]} & 10.0 & BOND & 1.00 & 0.33 \\
Distributed \texttt{[3.8,\ldots,3.8]} & 10.1 & FREEZE & 0.38 & 0.53 \\
\bottomrule
\end{tabular}
\end{table}

The pure vectors produce clear, high-activation actions with low conflict. The distributed vector, with the same energy spread across all axes, produces weak FREEZE with elevated conflict. This is the anxious-vigilance state: everything is mildly activated, nothing dominates, the organism is alert but undirected. A scalar ``arousal'' score would equate these four states. The valence vector distinguishes them.

The conflict score captures a phenomenon documented in behavioural neuroscience: decision latency increases when competing motivational circuits activate simultaneously \citep{mcnaughton2004}. The FREEZE response to \texttt{[fear=7, rage=7]} (conflict=0.10, paradoxically low because FREEZE itself dominates) differs from the distributed freeze (conflict=0.53). The former is tonic immobility, a coherent defensive strategy; the latter is paralysis by indecision.

\subsubsection{Competing Drives (Survival)}

The most interesting action states arise from axis interactions:

\begin{table}[h]
\centering
\caption{Axis interactions produce emergent actions.}
\label{tab:competing}
\begin{tabular}{llll}
\toprule
\textbf{Scenario} & \textbf{Vector} & \textbf{Action} & \textbf{Biological interpretation} \\
\midrule
Predator & \texttt{[0,0,8,0,0,0,0]} & FLEE (0.64) & Standard escape \\
Cornered & \texttt{[0,7,7,0,0,0,0]} & FREEZE (0.70) & Tonic immobility \\
Mother + predator & \texttt{[0,0,8,0,8,0,0]} & PROTECT/FLEE (tie) & Parental defence \\
Lost offspring & \texttt{[0,0,0,0,3,8,0]} & SEEK\_CONTACT (0.80) & Separation distress \\
\bottomrule
\end{tabular}
\end{table}

The PROTECT response to \texttt{[fear=8, care=8]} is biologically significant: it represents the well-documented phenomenon of maternal aggression, where a parent's care drive overrides fear to defend offspring \citep{lonstein2002}. Neither axis alone predicts this behaviour. It emerges from the interaction, from the direction of the vector in the fear-care plane.

\subsubsection{Trajectory Preemption (Fire)}

We simulate a warehouse fire deteriorating over 6 timesteps and evaluate the action field at each point, including the projected state at horizon $h=2$:

\begin{table}[h]
\centering
\caption{Action transitions and preemption signals in fire scenario.}
\label{tab:preemption}
\begin{tabular}{cllcc}
\toprule
\textbf{Time} & \textbf{Current action} & \textbf{Projected action} & \textbf{Preempt?} & \textbf{Conflict} \\
\midrule
$t=0$ & HOLD\_POSITION (0.45) & HOLD\_POSITION (0.45) & No & 1.05 \\
$t=1$ & HOLD\_POSITION (0.29) & DEFENSIVE (0.16) & No & 0.67 \\
$t=2$ & HOLD\_POSITION (0.10) & DEFENSIVE (0.30) & \textbf{Yes} & 1.14 \\
$t=3$ & EVACUATE (0.20) & EVACUATE (0.52) & No & 0.86 \\
$t=4$ & EVACUATE (0.42) & MAYDAY (0.94) & \textbf{Yes} & 0.20 \\
$t=5$ & EVACUATE (0.68) & MAYDAY (2.92) & \textbf{Yes} & 0.31 \\
\bottomrule
\end{tabular}
\end{table}

Three observations:

\begin{enumerate}
    \item \textbf{High conflict at transition points.} At $t=0$ and $t=2$, conflict exceeds 1.0 as multiple actions compete at near-equal strength. These are the moments when a firefighter's experience (calibrated valence heads) makes the difference between a good decision and a fatal one.

    \item \textbf{Preemption fires before the crisis.} At $t=2$, the current state still supports HOLD\_POSITION, but the projected state at $t+2$ shows DEFENSIVE dominating. The system signals that it is time to transition posture before conditions force it. At $t=4$, the current action is EVACUATE but the projection shows MAYDAY. The signal is unambiguous: leave now, because in 2 timesteps you will be calling for rescue.

    \item \textbf{Conflict drops as urgency rises.} By $t=4$ and $t=5$, conflict is low (0.20, 0.31) and the situation is unambiguous. High conflict is a feature of ambiguous, moderate-threat states. Clear emergencies produce clear actions. This matches the phenomenology of expert decision-making under stress: the hardest decisions are in the middle zone, not at the extremes \citep{klein1998}.
\end{enumerate}

\subsection{Bayesian Experience Memory}

A scoring system without memory is a thermometer. It measures but does not learn. We introduce a Bayesian experience memory that stores, retrieves, and selectively forgets experiences based on their statistical contribution to the system's posterior beliefs.

\subsubsection{The Learning Loop}

Each experience passes through a Bayesian loop:

\begin{enumerate}
    \item \textbf{Observe} visual pattern, extract features, score on domain axes
    \item \textbf{Update} the posterior: $P(\theta | D_{1:t}) \propto P(x_t | \theta) \cdot P(\theta | D_{1:t-1})$
    \item \textbf{Compute surprise}: $\mathrm{KL}[P(\theta | D_{1:t}) \| P(\theta | D_{1:t-1})]$
    \item \textbf{Gate}: store only if statistically informative (high surprise, extreme outcome, or novel pattern)
    \item \textbf{Evict} if at capacity, using posterior-aware selection
    \item \textbf{Replay} periodically: recompute all retention scores against the current posterior
\end{enumerate}

The posterior is parameterised as a diagonal Gaussian over valence outcomes: mean $\mu$ and precision $\lambda$ per axis, updated via conjugate Bayesian update:

\begin{equation}
\lambda_{\text{new}} = \lambda_{\text{old}} + \lambda_{\text{obs}}, \quad \mu_{\text{new}} = \frac{\lambda_{\text{old}} \mu_{\text{old}} + \lambda_{\text{obs}} x}{\lambda_{\text{new}}}
\end{equation}

Surprise is the KL divergence between old and new posterior, quantifying the amount of learning this observation caused.

\subsubsection{The Eviction Problem: Sequential Bias}

Naive eviction (remove lowest-scoring experience) introduces three compounding biases:

\begin{itemize}
    \item \textbf{Primacy bias}: early experiences receive inflated surprise because the prior was flat, a trivially true but uninformative artefact
    \item \textbf{Redundancy blindness}: 10 copies of ``routine fire'' are treated as 10$\times$ value when they are $\sim$1$\times$ value
    \item \textbf{Coverage drift}: extreme experiences dominate, routine ones are lost, and the stored distribution no longer represents the true distribution
\end{itemize}

We address these with three mechanisms:

\textbf{1. Normalised influence.} The leave-one-out (LOO) influence of experience $i$ is:

\begin{equation}
I_i = \mathrm{KL}[P(\theta | D) \| P(\theta | D \setminus \{x_i\})]
\end{equation}

This measures how much the posterior would change if $x_i$ were removed. However, raw $I_i$ is degenerate when $n$ is small (removing 1 of 2 observations trivially shifts the posterior). We normalise:

\begin{equation}
\hat{I}_i = \frac{I_i}{\sqrt{n_{\text{storage}}}}
\end{equation}

\noindent where $n_{\text{storage}}$ is the posterior size when $x_i$ was stored. This discounts early experiences whose high influence is an artefact of posterior thinness.

\begin{table}[h]
\centering
\caption{Effect of influence normalisation.}
\label{tab:influence}
\small
\begin{tabular}{llcccc}
\toprule
\textbf{Experience} & \textbf{Action} & \textbf{Raw influence} & $n_{\text{storage}}$ & \textbf{Normalised} & \textbf{Effect} \\
\midrule
\#0 (first ever) & MONITOR & 1.739 & 1 & 0.422 & Discounted: flat-prior artefact \\
\#8 (growing) & DEFENSIVE & 0.232 & 9 & 0.077 & Moderate discount \\
\#16 (crisis) & EVACUATE & 1.739 & 17 & 0.422 & Retains high score: genuinely informative \\
\#69 (routine) & MONITOR & 0.050 & 70 & 0.006 & Heavily discounted: posterior already sharp \\
\bottomrule
\end{tabular}
\end{table}

\textbf{2. Coverage-driven eviction.} Feature space is partitioned into regions via hyperoctant hashing. Eviction preferentially removes experiences from dense regions, rebalancing the stored distribution toward uniform coverage. A region with 1 experience is never emptied, maintaining representational coverage even under memory pressure.

\textbf{3. Importance-weighted replay (the backward pass).} Every $k$ insertions, all stored experiences are re-evaluated against the \emph{current} posterior:

\begin{itemize}
    \item \textbf{Recompute influence}: an experience that was load-bearing may now be redundant (later experiences covered the same territory)
    \item \textbf{Recompute surprise}: an experience that seemed routine may now be an outlier (similar experiences were evicted)
    \item \textbf{Importance weighting}: experiences far from the current posterior mean receive amplified surprise because they are informative precisely when they are unusual
\end{itemize}

This replay is the backward pass: propagating the current posterior state back through all stored experiences and recomputing their value. Without it, retention scores become stale and the eviction policy degrades.

\subsubsection{Results: What Memory Keeps}

We simulate 80 fire observations (61 routine small fires, 10 growing fires, 3 crises, 2 backdrafts, 1 electrical fire) through a memory with capacity 20.

\begin{table}[h]
\centering
\caption{Memory composition after 80 observations.}
\label{tab:memorycomp}
\begin{tabular}{lccc}
\toprule
\textbf{Profile} & \textbf{Seen} & \textbf{In memory} & \textbf{Retention rate} \\
\midrule
Crisis/Backdraft/Electrical & 6 & 6 & 100\% \\
Growing fire & 10 & 5 & 50\% \\
Routine small & 61 & 9 & 15\% \\
\bottomrule
\end{tabular}
\end{table}

All rare, high-consequence events are retained. Growing fires (moderate intensity) retain representatives. Routine fires are heavily pruned because the posterior already knows what they look like.

\textbf{Temporal bias check:}

\begin{table}[h]
\centering
\begin{tabular}{lcc}
\toprule
\textbf{Period} & \textbf{Retained} & \textbf{Expected (unbiased)} \\
\midrule
Early (first 20 obs) & 9 & 5.2 \\
Middle (obs 20--49) & 7 & 7.8 \\
Late (obs 50+) & 4 & 7.0 \\
\bottomrule
\end{tabular}
\end{table}

The middle period, where the growing fires and first crisis occurred, is well-represented (7 vs.\ expected 7.8). Early-period overrepresentation (9 vs.\ 5.2) persists but is reduced 40\% compared to naive eviction (which retained 15/20 from the early period). The normalised influence correction is working but could be strengthened with a more aggressive discount function.

\subsubsection{Connection to Biological Memory}

The three mechanisms mirror known properties of biological memory consolidation:

\begin{itemize}
    \item \textbf{Influence $\leftrightarrow$ synaptic tagging}: hippocampal memories are tagged for consolidation based on their novelty relative to existing schema \citep{wang2010}
    \item \textbf{Coverage $\leftrightarrow$ pattern separation}: the dentate gyrus actively separates similar inputs to maintain distinct representations \citep{leutgeb2007}
    \item \textbf{Replay $\leftrightarrow$ sleep consolidation}: offline replay during sleep reactivates memories and re-evaluates them against updated cortical representations \citep{diekelmann2010}
\end{itemize}

The parallel is not metaphorical. Biological memory systems face the same statistical problem: finite storage, non-stationary input distribution, and the need to retain informative experiences while discarding redundant ones. The solution, posterior-aware selective retention with periodic replay, converges from both the computational and biological directions.

\section{Discussion}

\subsection{The Hierarchy Inverted}

Current vision-language models place language at the centre: visual features are projected into a text embedding space, evaluated by text-based loss functions, and benchmarked by text generation quality. vScore inverts this hierarchy:

\begin{table}[h]
\centering
\begin{tabular}{lll}
\toprule
\textbf{Layer} & \textbf{Vision-Language Models} & \textbf{vScore} \\
\midrule
Foundation & Language tokens & Valence scores \\
Training signal & Text supervision & Outcome scores \\
Output & Words & Trigger events \\
Language role & Core & Optional serialisation \\
\bottomrule
\end{tabular}
\end{table}

This inversion is not merely architectural. It changes what the system optimises for. A VLM optimises for describing what it sees. vScore optimises for predicting what will happen and whether it matters. The former produces narration. The latter produces intelligence.

\subsection{Universal Primitives vs.\ Domain-Specific Axes}

Our transfer experiments reveal a natural partition of visual features:

\textbf{Universal primitives} (transfer across domains):
\begin{itemize}
    \item Motion magnitude (speed)
    \item Object trajectory (direction, acceleration)
    \item Spatial extent (how much of the visual field is activated)
\end{itemize}

\textbf{Domain-specific features} (require targeted training):
\begin{itemize}
    \item Impact dynamics (collision patterns)
    \item Coordination patterns (synchronised multi-agent motion)
    \item Tension dynamics (stillness-before-release)
    \item Domain-specific textures (flame, ice, water)
\end{itemize}

This partition aligns with neuroscience: early visual processing (V1, MT) extracts motion, edges, and spatial frequency as universal primitives, while higher areas (IT, STS) encode category-specific and action-specific representations. vScore's architecture mirrors this: the frozen encoder provides universal primitives; the trainable valence heads learn domain-specific scoring.

Critically, the 13-category experiment revealed that the partition itself is data-dependent: coordination appeared domain-specific with 5 categories but universal with 13. This suggests that the boundary between universal and domain-specific is not fixed; it shifts as the training set diversifies. More extensive research with larger, more diverse datasets is required to establish the definitive partition.

\subsection{From Detection to Understanding: The Geometry Argument}

Motion detection is trivially solved. Any optical flow algorithm detects motion. The question is what motion \emph{means}. vScore's answer is that meaning is not a label attached to motion but a \emph{direction in valence space} that the motion contributes to.

Consider the transition from detection to understanding:

\begin{enumerate}
    \item \textbf{Detection}: ``something is moving'' (optical flow, background subtraction, solved in the 1980s)
    \item \textbf{Scoring}: ``it is moving at speed=7 toward proximity=8'' (per-axis regression, Section~4.3)
    \item \textbf{Geometry}: ``the valence vector is pointing toward the EVACUATE region of action space'' (Section~4.5)
    \item \textbf{Projection}: ``and the trajectory is curving toward MAYDAY'' (Section~4.5.3)
\end{enumerate}

Each level adds information that the previous level lacks. Detection without scoring is noise. Scoring without geometry is a dashboard of unrelated numbers. Geometry without projection is a snapshot without consequences.

The action field formalisation makes this concrete: an $n$-dimensional valence space contains a continuous manifold of behavioural states. Each point is characterised not by which axis is highest (that would reduce to classification) but by its position relative to the boundaries between action regions. The boundaries themselves are nonlinear surfaces in valence space. For example, the FREEZE region lies along the diagonal where fear and rage are approximately equal, not at any single threshold.

This is fundamentally different from both classification (``this is a fire'') and regression (``the fire is at intensity 7.5''). It is \emph{geometric inference}: given where the vector is, where it is heading, and the shape of the action landscape, what response does the geometry demand?

\subsection{The Naming Problem: Words Imposed on Wordless Representations}

There is an unavoidable irony in this work. We argue that intelligence operates below language, then we name the axes ``speed,'' ``impact,'' ``tension.'' These names are our linguistic projection onto what the encoder is actually doing, and every name is a lossy compression.

The encoder has no concept of speed in the physical sense of distance over time. It has no ruler, no clock, no scale. What it detects is the rate of spatial displacement of correlated feature clusters between frames. Pixels that were here are now there, and the difference between frames captures that displacement. We call this ``speed'' because it is the closest word available to us. A pre-linguistic organism does not call it anything. It reacts to the rate of change without naming it.

This matters because the six axes we defined (speed, impact, precision, verticality, coordination, tension) are not ground truth. They are our best available approximation, imposed from outside, of what the encoder's 1024 dimensions naturally represent. The real structure of V-JEPA~2's representation space may contain axes that do not correspond to any single word, or axes that blend what we would separate linguistically. When we train a valence head to predict ``speed = 7,'' we are asking the encoder to collapse its rich internal representation onto a single human-legible number. Information is inevitably lost in that projection.

A stronger version of this framework would let the axes emerge from the data rather than defining them in advance. Unsupervised clustering of the encoder's representation space, followed by post-hoc interpretation of what each cluster corresponds to, would reveal the encoder's own ontology rather than imposing ours. The named axes served their purpose for this initial work: they made the framework testable and the results interpretable. But the named axes are scaffolding. The building is the representation space itself.

We are, in the end, forced to use words to describe a system that operates without them. This limitation belongs to us, not to the system.

\subsection{Implications for Annotation}

Traditional annotation asks: ``What is in this video?'' (a linguistic question). vScore annotation asks: ``How activated is each axis at this moment?'' (a numerical question).

This has practical advantages:

\begin{enumerate}
    \item \textbf{No linguistic competence required.} Annotators can score videos by adjusting sliders, regardless of language or literacy.
    \item \textbf{Cross-cultural validity.} A Japanese firefighter and a Brazilian firefighter may use different words, but they read the same visual dynamics and would produce similar valence scores.
    \item \textbf{Physiological validation.} Valence scores can be validated against heart rate, galvanic skin response, and eye tracking, providing ground truth that is independent of language.
    \item \textbf{Continuous, not categorical.} ``How fast?'' (0--10) captures more information than ``fast/slow.''
\end{enumerate}

\subsection{Toward Predictive Visual Intelligence}

The trajectory projection and threshold triggering components of vScore formalise a capacity that biological systems exhibit but current AI systems largely lack: \textbf{acting on projected futures, not observed presents.}

A self-driving car that detects a pedestrian after they step into the road is reactive. One that detects the trajectory of a person walking toward the curb and projects the crossing is predictive. The difference is not better object detection. It is valence trajectory projection.

vScore provides the computational framework for this: score the current state, track the trajectory, project forward, trigger when the projection crosses a threshold. The visual encoder is a commodity (V-JEPA~2 today, its successor tomorrow). The intelligence is in the scoring and projection.

\subsection{Limitations}

\begin{enumerate}
    \item \textbf{Proxy annotations.} Our current experiments use fixed per-category scores. Real per-video annotations from domain experts or physiological measurements would enable learning intra-category variation, for instance two archery videos with different tension profiles.

    \item \textbf{Small dataset.} 50 videos across 5 categories. The transfer findings are directional, not statistically conclusive. Validation on larger, more diverse datasets (Kinetics-700, Ego4D, domain-specific collections) is needed.

    \item \textbf{Global pooling.} We average 8{,}192 spatiotemporal tokens to a single vector, discarding spatial and temporal structure. The dense token representation from V-JEPA~2 is designed for spatial grounding, and future work should score per-region and per-timestep, not per-video.

    \item \textbf{Static thresholds.} The current threshold mechanism uses fixed base values. In biological systems, thresholds adapt to context (heightened alertness, habituation, priming). Learned, context-dependent thresholds are a natural extension.

    \item \textbf{Single encoder.} We test only V-JEPA~2. The framework is encoder-agnostic, and testing with alternative self-supervised encoders (DINOv2, VideoMAE, InternVideo) would isolate the contribution of the scoring framework from the encoder quality.
\end{enumerate}

\section{Related Work}

\textbf{Joint Embedding Predictive Architectures.} \citeauthor{lecun2022}'s (\citeyear{lecun2022}) position paper on autonomous machine intelligence argued that autoregressive generative models, including LLMs, are insufficient for world understanding because they operate on discrete tokens rather than continuous representations. He proposed JEPA as an alternative: learn to predict representations of the world from other representations, without reconstructing inputs. This is a fundamental departure from both generative models (which predict pixels or tokens) and contrastive learning (which only learns to distinguish). The key insight for our work is that JEPA representations are, by design, \emph{not grounded in language}. They encode visual structure as the world presents it, not as language describes it. This makes them the ideal substrate for pre-linguistic valence scoring. I-JEPA \citep{assran2023} demonstrated this principle for images; V-JEPA \citep{bardes2024} and V-JEPA~2 \citep{bardes2025} extended it to video, learning dense spatiotemporal features through masked prediction in latent space. V-JEPA~2 in particular produces the kind of spatially structured, temporally consistent representations that valence scoring requires: features that track objects, motion, and spatial relationships across time without any text supervision. vScore can be understood as a downstream answer to the question LeCun's architecture poses: if the encoder sees the world without words, what should the evaluation layer look like?

\textbf{Self-supervised video representation learning.} Beyond the JEPA family, self-supervised video learning includes contrastive approaches \citep{qian2021}, masked autoencoders \citep{tong2022,wang2023}, and distillation methods \citep{oquab2023}. vScore is encoder-agnostic; any self-supervised video encoder that produces dense features could serve as Level~2. We chose V-JEPA~2 because its JEPA-based training most closely aligns with the principle that visual understanding should be independent of language.

\textbf{Affective computing.} Emotion recognition from video \citep{mollahosseini2017,li2020} typically maps visual input to categorical emotion labels (happy, sad, angry) or dimensional affect (valence-arousal). vScore differs in two ways: (1)~the axes are domain-specific and task-relevant, not universal emotion categories, and (2)~the emphasis is on trajectory projection and threshold triggering, not state classification.

\textbf{Affordance detection.} Work on visual affordances \citep{nagarajan2019,luo2022} detects action possibilities in images. vScore extends this from ``what can be done'' to ``what will happen and how much does it matter,'' adding temporal projection and outcome valence.

\textbf{Predictive coding.} Friston's Free Energy Principle \citep{friston2010} frames perception as prediction error minimisation. vScore operationalises a related idea: the system projects expected trajectories and triggers on deviations from expected homeostasis. The analogy is direct: surprise (prediction error) corresponds to non-zero valence magnitude.

\textbf{Embodied AI and robotics.} V-JEPA~2 reports a 20-point improvement in robotic grasping success, demonstrating that dense visual features support motor control. vScore provides a framework for the evaluation layer between perception and action, addressing the ``should I act?'' decision that precedes ``how do I act?''

\section{Future Directions}

\subsection{Cross-Modal Valence (Preliminary Results)}

The vScore mechanism is modality-agnostic by construction. The valence vector is the universal interface. Downstream of it, nothing knows whether the input was pixels or pressure waves. We formalise this with a multimodal architecture where each modality has its own encoder and valence head, but all heads output to the same $N$-axes valence space. Fusion happens in valence space, not feature space.

This design makes modality conflict observable and diagnostic. In a simulated ``occluded threat'' scenario (growl heard but source not visible), the audio head scores fear=8 while the visual head scores fear=1, producing a conflict signal of 4.9 on the fear axis. This conflict itself is informative: audio threat $\gg$ visual threat implies the danger source is occluded or approaching from outside the visual field, a signal that would be lost in feature-space fusion.

We define three fusion rules: MAX (survival default: if any sense says danger, respond), MEAN (requires cross-modal confirmation), and confidence-weighted (trust the modality with more experience). The fusion rule is a domain parameter, not a fixed architectural choice. A survival domain uses MAX. A diagnostic domain like industrial machine monitoring might use CONFLICT to detect when auditory and visual signals diverge.

We implement three audio domains: survival sound (threat, proximity, urgency, familiarity, social signal, rhythmic pull, dissonance), industrial sound (anomaly, degradation, impact event, resonance shift, load stress), and music (tension, release, energy, intimacy, novelty, groove, melancholy). These demonstrate that the same metaclass architecture applies without modification. The principle holds: a baby covers its ears before it knows the word ``loud.'' A dog runs from thunder without vocabulary. The evaluation is pre-linguistic and pre-modal.

\subsection{Learned Threshold Adaptation}

Current thresholds are fixed. Biological thresholds adapt: a soldier in combat has lower FEAR thresholds (hypervigilance); a child at play has higher ones (safety). Learning context-dependent threshold functions from data is a natural next step.

\subsection{Compositional Domains}

Can a system trained on Fire and Weather generalise to ``wildfire during a storm''? Compositional domain scoring, where multi-domain valence vectors interact, would test whether the framework supports the combinatorial richness of real-world scenarios.

\subsection{Expert Calibration Studies}

The strongest validation would be: have domain experts (firefighters, coaches, traders) score videos on their domain axes, train valence heads on their scores, and test whether the resulting model predicts expert attention and decision-making. If vScore heads correlate with expert gaze patterns and response times, the framework captures something real about expert visual intelligence.

\section{Conclusion}

We have presented vScore, a framework that reframes visual AI from ``what is this?'' to ``what does this mean for outcome-relevant axes?'' The core claim is that visual intelligence, in biological systems and potentially in artificial ones, operates pre-linguistically through valence scoring, trajectory projection, and threshold triggering. Language is a late-stage serialisation, not a prerequisite for understanding.

Our experiments across 213 videos from 13 categories and two datasets demonstrate that self-supervised video features contain universal visual primitives, specifically coordination, impact, and speed, that transfer to unseen domains with low error (MAE 1.45 to 1.74). The 13-category experiment revealed that the boundary between universal and domain-specific is itself data-dependent: coordination appeared domain-specific with 5 categories but universal with 13, suggesting that the encoder already captures these dynamics but they simply require sufficient diversity of training examples to be extracted. Tension remains genuinely domain-specific (MAE 3.13, transfers in only 38\% of holdouts), confirming that some visual dynamics are irreducibly tied to their context.

These results are promising but preliminary. More extensive research is required to establish the definitive partition of universal and domain-specific visual primitives, to validate the framework with real expert annotations rather than proxy scores, to integrate audio encoders for cross-modal valence scoring, and to test the Bayesian memory system on temporal sequences of real-world video. The framework's strength lies in its simplicity and extensibility. A scored domain is defined by a list of axes, a valence head is a small MLP, and the entire downstream pipeline (action geometry, trajectory projection, threshold triggering, Bayesian memory) operates identically regardless of domain or modality. The approach shows promise as a foundation for pre-linguistic visual intelligence; the work ahead is to scale it.

Zero is safety. Everything above zero is cost. The system exists to return to zero.

\section*{Acknowledgements}

This work owes a foundational debt to Yann LeCun, whose articulation of the Joint Embedding Predictive Architecture \citep{lecun2022} provided both the theoretical grounding and the practical encoder on which vScore is built. LeCun's central argument, that autonomous intelligence requires world models predicting in representation space rather than token space and that autoregressive language models are fundamentally limited by their confinement to discrete symbols, is the premise from which vScore departs. Without V-JEPA and V-JEPA~2, developed by LeCun's team at Meta FAIR (Bardes, Assran, Ballas, Rabbat, LeCun, and collaborators), the empirical results in this paper would not exist. The idea that a visual encoder can learn dense, temporally consistent representations without any language supervision is not our contribution; it is theirs. Our contribution is to ask what comes \emph{after} the encoder when language is deliberately excluded from the evaluation pipeline.

We also acknowledge the broader JEPA research programme at Meta FAIR, including I-JEPA \citep{assran2023}, whose demonstration that joint embedding prediction produces semantically meaningful image features without pixel reconstruction or text alignment was an essential precursor to the video-domain extensions we rely on.

\bibliographystyle{plainnat}

\begin{thebibliography}{99}

\bibitem[Assran et~al.(2023)]{assran2023}
Assran, M., et~al. (2023).
\newblock Self-Supervised Learning from Images with a Joint-Embedding Predictive Architecture.
\newblock \emph{CVPR}.

\bibitem[Bardes et~al.(2024)]{bardes2024}
Bardes, A., et~al. (2024).
\newblock V-JEPA: Latent Video Prediction for Visual Representation Learning.
\newblock \emph{arXiv preprint}.

\bibitem[Bardes et~al.(2025)]{bardes2025}
Bardes, A., et~al. (2025).
\newblock V-JEPA 2: Unlocking Dense Features in Video Self-Supervised Learning.
\newblock \emph{arXiv:2603.14482}.

\bibitem[Corr(2013)]{corr2013}
Corr, P.~J. (2013).
\newblock Approach and Avoidance Behaviour: Multiple Systems and Their Interactions.
\newblock \emph{Emotion Review}, 5(3), 285--290.

\bibitem[Damasio(1994)]{damasio1994}
Damasio, A. (1994).
\newblock \emph{Descartes' Error: Emotion, Reason, and the Human Brain}.
\newblock Putnam.

\bibitem[Diekelmann \& Born(2010)]{diekelmann2010}
Diekelmann, S., \& Born, J. (2010).
\newblock The Memory Function of Sleep.
\newblock \emph{Nature Reviews Neuroscience}, 11(2), 114--126.

\bibitem[Friston(2010)]{friston2010}
Friston, K. (2010).
\newblock The Free-Energy Principle: A Unified Brain Theory?
\newblock \emph{Nature Reviews Neuroscience}, 11(2), 127--138.

\bibitem[Gallup(1977)]{gallup1977}
Gallup, G.~G. (1977).
\newblock Tonic Immobility: The Role of Fear and Predation.
\newblock \emph{The Psychological Record}, 27(1), 41--61.

\bibitem[Gibson(1979)]{gibson1979}
Gibson, J.~J. (1979).
\newblock \emph{The Ecological Approach to Visual Perception}.
\newblock Houghton Mifflin.

\bibitem[Gourgari et~al.(2013)]{gourgari2013}
Gourgari, S., et~al. (2013).
\newblock THETIS: Three Dimensional Tennis Shots, a Human Action Dataset.
\newblock \emph{CVPR Workshops}.

\bibitem[Klein(1998)]{klein1998}
Klein, G. (1998).
\newblock \emph{Sources of Power: How People Make Decisions}.
\newblock MIT Press.

\bibitem[LeCun(2022)]{lecun2022}
LeCun, Y. (2022).
\newblock A Path Towards Autonomous Machine Intelligence.
\newblock \emph{OpenReview preprint}, Version 0.9.2.

\bibitem[Leutgeb et~al.(2007)]{leutgeb2007}
Leutgeb, J.~K., et~al. (2007).
\newblock Pattern Separation in the Dentate Gyrus and CA3 of the Hippocampus.
\newblock \emph{Science}, 315(5814), 961--966.

\bibitem[Li \& Deng(2020)]{li2020}
Li, S., \& Deng, W. (2020).
\newblock Deep Facial Expression Recognition: A Survey.
\newblock \emph{IEEE Transactions on Affective Computing}.

\bibitem[Lonstein \& Gammie(2002)]{lonstein2002}
Lonstein, J.~S., \& Gammie, S.~C. (2002).
\newblock Sensory, Hormonal, and Neural Control of Maternal Aggression in Laboratory Rodents.
\newblock \emph{Neuroscience \& Biobehavioural Reviews}, 26(8), 869--888.

\bibitem[Luo et~al.(2022)]{luo2022}
Luo, H., et~al. (2022).
\newblock Learning Affordance Grounding from Exocentric Images.
\newblock \emph{CVPR}.

\bibitem[McNaughton \& Corr(2004)]{mcnaughton2004}
McNaughton, N., \& Corr, P.~J. (2004).
\newblock A Two-Dimensional Neuropsychology of Defence: Fear/Anxiety and Defensive Distance.
\newblock \emph{Neuroscience \& Biobehavioural Reviews}, 28(3), 285--305.

\bibitem[Mollahosseini et~al.(2017)]{mollahosseini2017}
Mollahosseini, A., et~al. (2017).
\newblock AffectNet: A Database for Facial Expression, Valence, and Arousal Computing in the Wild.
\newblock \emph{IEEE Transactions on Affective Computing}.

\bibitem[Nagarajan et~al.(2019)]{nagarajan2019}
Nagarajan, T., et~al. (2019).
\newblock Grounded Human-Object Interaction Hotspots from Video.
\newblock \emph{ICCV}.

\bibitem[Oquab et~al.(2023)]{oquab2023}
Oquab, M., et~al. (2023).
\newblock DINOv2: Learning Robust Visual Features without Supervision.
\newblock \emph{arXiv preprint}.

\bibitem[Panksepp(1998)]{panksepp1998}
Panksepp, J. (1998).
\newblock \emph{Affective Neuroscience: The Foundations of Human and Animal Emotions}.
\newblock Oxford University Press.

\bibitem[Qian et~al.(2021)]{qian2021}
Qian, R., et~al. (2021).
\newblock Spatiotemporal Contrastive Video Representation Learning.
\newblock \emph{CVPR}.

\bibitem[Regan \& Vincent(1995)]{regan1995}
Regan, D., \& Vincent, A. (1995).
\newblock Visual Processing of Looming and Time to Contact Throughout the Visual Field.
\newblock \emph{Vision Research}, 35(13), 1845--1857.

\bibitem[Schiff(1965)]{schiff1965}
Schiff, W. (1965).
\newblock Perception of Impending Collision: A Study of Visually Directed Avoidant Behaviour.
\newblock \emph{Psychological Monographs}, 79(11), 1--26.

\bibitem[Tong et~al.(2022)]{tong2022}
Tong, Z., et~al. (2022).
\newblock VideoMAE: Masked Autoencoders are Data-Efficient Learners for Self-Supervised Video Pre-Training.
\newblock \emph{NeurIPS}.

\bibitem[Wang \& Morris(2010)]{wang2010}
Wang, S.-H., \& Morris, R.~G.~M. (2010).
\newblock Hippocampal-Neocortical Interactions in Memory Formation, Consolidation, and Reconsolidation.
\newblock \emph{Annual Review of Psychology}, 61, 49--79.

\bibitem[Wang et~al.(2023)]{wang2023}
Wang, L., et~al. (2023).
\newblock VideoMAE V2: Scaling Video Masked Autoencoders with Dual Masking.
\newblock \emph{CVPR}.

\bibitem[Zhao et~al.(2019)]{zhao2019}
Zhao, H., et~al. (2019).
\newblock HACS: Human Action Clips and Segments Dataset for Recognition and Temporal Localization.
\newblock \emph{ICCV}.

\end{thebibliography}

\appendix

\section{Project Structure}

\begin{lstlisting}[language={}]
vScore/
|-- core/
|   |-- metaclass.py      # ScoredDomain metaclass
|   |-- valence.py         # ValenceVector, ValenceTrajectory
|   +-- threshold.py       # ThresholdTrigger, dynamic thresholds
|-- domains/
|   |-- survival.py        # Panksepp's 7 circuits
|   |-- fire.py            # 7 axes
|   |-- weather.py         # 6 axes
|   |-- hockey.py          # 6 axes
|   |-- trading.py         # 6 axes
|   |-- sound.py           # 7 axes, pre-linguistic auditory scoring
|   |-- industrial.py      # 6 axes, machine/environment audio
|   +-- music.py           # 7 axes, affective musical scoring
|-- encoder/
|   +-- bridge.py          # V-JEPA 2 -> valence head bridge
|-- projection/
|   +-- temporal.py        # Transformer-based trajectory projector
|-- data/
|   +-- schema.py          # Annotation format (no words, just scores)
|-- train.py               # Training pipeline
|-- run_cross_domain.py    # Cross-domain feature analysis
|-- demo_actions.py        # Action inference from valence geometry
|-- demo_memory.py         # Naive Bayesian memory
|-- demo_memory_bayesian.py # Posterior-aware eviction
+-- demo_multimodal.py     # Cross-modal fusion in valence space
\end{lstlisting}

Code available at: \url{https://github.com/Tennisee-data/vScore}

\section{Reproducibility}

All experiments run on Apple M-series CPU (no GPU required for inference and training at this scale). Feature extraction for 213 videos takes $\sim$30 minutes (one-time; features are cached as 1{,}024-dim vectors, 1.3\,MB total). Training each valence head takes $<$5 seconds (500 epochs on cached vectors). Full cross-domain transfer experiment (13 holdout runs) takes $<$60 seconds. Total reproduction time from scratch: $\sim$35 minutes.

Dependencies: PyTorch 2.8, Transformers 4.57, PyAV 16.1, yt-dlp (for HACS download), V-JEPA~2 weights from \texttt{facebook/vjepa2-vitl-fpc64-256}.

Data sources: Kinetics-mini (\texttt{nateraw/kinetics-mini} on HuggingFace), HACS \citep{zhao2019}; \texttt{github.com/hangzhaomit/HACS-dataset}.

\end{document}
